% Standalone XeLaTeX source; source problem: 2014ap_exam.pdf, PDF pp. 10--11.
\documentclass[11pt,a4paper]{article}
\usepackage[margin=24mm]{geometry}
\usepackage{fontspec}
\usepackage{xeCJK}
\setCJKmainfont{HaranoAjiMincho-Regular.otf}
\setCJKsansfont{HaranoAjiGothic-Medium.otf}
\usepackage{amsmath,amssymb}
\usepackage{tikz}
\usetikzlibrary{arrows.meta}
\setlength{\parindent}{1em}
\setlength{\parskip}{0.35em}
\title{物理工学専攻 院試解答集 ver2\\2014年 物理学II 第2問}
\author{政岡凜太郎}
\date{}
\begin{document}
\maketitle

$\beta=1/(k_BT)$ とする。粒子は古典Boltzmann統計に従い、$N!$ で割って同種粒子の重複計数を除く。

\section*{[1] 理想気体の分配関数}
一粒子の並進分配関数は位置積分と三つの運動量積分の積であり
\begin{align}
 z_1&=\frac{1}{h^3}\int_Vd^3r\int_{\mathbb R^3}d^3p\,
 e^{-\beta p^2/(2m)}\\
 &=\frac{V}{h^3}\left(\int_{-\infty}^{\infty}
 e^{-\beta p_x^2/(2m)}dp_x\right)^3
 =V\left(\frac{mk_BT}{2\pi\hbar^2}\right)^{3/2}.
\end{align}
ここで $h=2\pi\hbar$ とガウス積分を用いた。従って
\begin{equation}
 \boxed{Z_0(T,V,N)=\frac{z_1^N}{N!}
 =\frac{V^N}{N!}\left(\frac{mk_BT}{2\pi\hbar^2}\right)^{3N/2}.}
\end{equation}

\section*{[2] 内部エネルギーと定積比熱}
\begin{equation}
 U=-\frac{\partial}{\partial\beta}\ln Z_0
 =\boxed{\frac32Nk_BT},\qquad
 C_V=\left(\frac{\partial U}{\partial T}\right)_{V,N}
 =\boxed{\frac32Nk_B}.
\end{equation}

\section*{[3] 理想気体の状態方程式}
$F_0=-k_BT\ln Z_0$ であり、$V$ に依存する項は $-Nk_BT\ln V$ だけなので
\begin{equation}
 P=-\left(\frac{\partial F_0}{\partial V}\right)_{T,N}
 =k_BT\frac{\partial\ln Z_0}{\partial V}
 =\frac{Nk_BT}{V}.
\end{equation}
従って
\begin{equation}
 \boxed{PV=Nk_BT.}
\end{equation}

\section*{[4] 引力と排除体積}
原問の一粒子エネルギーは $p^2/(2m)-\alpha N/V$ である。従って全ポテンシャルエネルギーは $-\alpha N^2/V$、Boltzmann因子は $\exp(\beta\alpha N^2/V)$ となる。短距離斥力は指定どおり配置空間体積を $V-Nb$ に置換するから
\begin{equation}
 \boxed{Z(T,V,N)=\frac{(V-Nb)^N}{N!}
 \left(\frac{mk_BT}{2\pi\hbar^2}\right)^{3N/2}
 \exp\!\left(\frac{\alpha N^2}{Vk_BT}\right).}
\end{equation}
Helmholtz自由エネルギーは
\begin{equation}
 F=-k_BT\ln Z=F_{\rm kin}(T,N)-Nk_BT\ln(V-Nb)
 -\frac{\alpha N^2}{V},
\end{equation}
ここで $F_{\rm kin}$ は $V$ に依存しない。よって
\begin{equation}
 \boxed{P=-F_V=\frac{Nk_BT}{V-Nb}-\frac{\alpha N^2}{V^2}},\qquad
 \boxed{\left(P+\frac{\alpha N^2}{V^2}\right)(V-Nb)=Nk_BT.}
\end{equation}

\section*{[5] 自由エネルギー最小と安定性}
固定した $T,N,P$ に対して $A(V)=F(T,N,V)+PV$ とする。停留点では
\begin{equation}
 A'(V)=F_V+P=0\quad\Longleftrightarrow\quad P=-F_V,
\end{equation}
右辺の $P(V)=-F_V$ は均一状態の圧力であり、固定した外圧 $P$ と停留点で一致する。さらに
\begin{equation}
 A''(V)=F_{VV}=-\left(\frac{\partial P}{\partial V}\right)_{T,N}.
\end{equation}
局所安定な熱平衡は $A''>0$ の極小なので $\partial P/\partial V<0$。$A''<0$ の極大は不安定である。$A''=0$ では二階微分だけで安定性を判定できず、この等温線で安定枝が終わる点は限界点となる。図1では左右の等しい最小値が同じGibbs自由エネルギーを持つ共存状態、中央の極大が不安定状態である。

\begin{center}
\begin{tikzpicture}[>=Stealth,x=0.32cm,y=120cm]
 \pgfmathsetmacro{\pc}{0.018081606931683265}
 \pgfmathsetmacro{\gzero}{-0.4685322546678666}
 \draw[->] (0,0)--(15,0) node[right] {$V/(Nb)$};
 \draw[->] (0,0)--(0,0.032) node[above] {$(F+PV)/(N\alpha/b)$};
\draw[thick] plot[domain=1.3:14,samples=241] (\x,{-0.25*ln(\x-1)-1/\x+\pc*\x-\gzero});
 \fill (1.6447993625905433,0) circle (1.3pt) node[below] {安定};
 \fill (9.723323053742876,0) circle (1.3pt) node[below] {安定};
 \fill (3.458082353218395,0.0170369489645495) circle (1.3pt) node[above] {不安定};
\end{tikzpicture}
\end{center}

図では $k_BT b/\alpha=1/4$ の例を用いた。共存圧力は
$P_{\rm coex}b^2/\alpha\simeq0.01808161$ であり、二つの極小の高さが等しくなるように求めた。
自由エネルギー図の縦軸は、体積によらない項と共存時の極小値を差し引いて示した。

\section*{[6] 圧力による傾きの変化}
圧力を $P+\delta P$ に変えると
\begin{equation}
 A_{P+\delta P}(V)=A_P(V)+\delta P\,V.
\end{equation}
したがって $\delta P>0$ では大体積側ほど自由エネルギーが余分に上がり、小体積側の極小がより低くなる。圧力を下げると逆に大体積側の極小が選ばれる。これは外圧を上げると圧縮が進み、外圧を下げると膨張・気化が進むことに対応する。
\begin{center}
\begin{tikzpicture}[>=Stealth,x=0.32cm,y=90cm]
 \pgfmathsetmacro{\pc}{0.018081606931683265}
 \pgfmathsetmacro{\gzero}{-0.4685322546678666}
 \draw[->] (0,-0.01)--(15,-0.01) node[right] {$V/(Nb)$};
 \draw[->] (0,-0.01)--(0,0.04) node[above] {$(F+PV)/(N\alpha/b)$};
 \draw[dashed] plot[domain=1.3:14,samples=241] (\x,{-0.25*ln(\x-1)-1/\x+\pc*\x-\gzero});
 \draw[thick] plot[domain=1.3:14,samples=241] (\x,{-0.25*ln(\x-1)-1/\x+(\pc+0.001)*\x-\gzero});
 \draw[dotted,thick] plot[domain=1.3:14,samples=241] (\x,{-0.25*ln(\x-1)-1/\x+(\pc-0.001)*\x-\gzero});
\end{tikzpicture}
\end{center}
破線は共存圧力、実線は圧力増加、点線は圧力減少を表す。圧力変化は $\delta P b^2/\alpha=\pm0.001$ とし、同じ自由エネルギーに $\delta P V$ を加えて描いた。

\section*{[7] 状態方程式と平衡等温線}
状態方程式 $P(V)=Nk_BT/(V-Nb)-\alpha N^2/V^2$ をこの温度で描くと、図1に二つの極小がある条件ではS字状のvan der Waalsループを持つ。
中間枝は $\partial P/\partial V>0$ であり、均一な状態として不安定である。
共存圧力 $P_{\rm coex}$ では、小体積相と大体積相のGibbs自由エネルギーが等しい。
全 $N$ 個の粒子がそれぞれの均一相にあるときの体積を $V_l,V_g$ と書き、$V_l<V_g$ とする。
$F(V)=F(T,N,V)$ と略記すると、二つの大域最小の高さが等しい条件は
\begin{equation}
 F(V_g)-F(V_l)=-P_{\rm coex}(V_g-V_l)
\end{equation}
である。

中間の全体積では、二つの相が共存し、各相に属する粒子の割合が変わる。
熱力学極限で界面の自由エネルギーを無視し、各相の体積と自由エネルギーを加える。
大体積相に属する粒子の割合を $x$（$0\le x\le1$）とすると、自由エネルギーの示量性から
\begin{align}
 V&=(1-x)V_l+xV_g,\\
 F_{\rm mix}(V)&=(1-x)F(V_l)+xF(V_g).
\end{align}
従ってこの混合状態の圧力は
\begin{equation}
 -\frac{dF_{\rm mix}}{dV}
 =-\frac{F(V_g)-F(V_l)}{V_g-V_l}=P_{\rm coex}.
\end{equation}
このため平衡等温線は $V_l\le V\le V_g$ で水平となる。
均一状態のS字曲線と、二相の割合が変わる平衡状態を同じ図に示す。
\begin{center}
\begin{tikzpicture}[>=Stealth,x=0.32cm,y=60cm]
 \pgfmathsetmacro{\pc}{0.018081606931683265}
 \draw[->] (0,0)--(15,0) node[right] {$V/(Nb)$};
 \draw[->] (0,0)--(0,0.065) node[above] {$Pb^2/\alpha$};
 \draw[dashed] plot[domain=1.5:14,samples=301] (\x,{0.25/(\x-1)-1/(\x*\x)});
 \draw[thick] plot[domain=1.5:1.6447993625905433,samples=41] (\x,{0.25/(\x-1)-1/(\x*\x)});
 \draw[thick] (1.6447993625905433,\pc)--(9.723323053742876,\pc);
 \draw[thick] plot[domain=9.723323053742876:14,samples=81] (\x,{0.25/(\x-1)-1/(\x*\x)});
 \node[above] at (7,\pc) {$P_{\rm coex}b^2/\alpha$};
\end{tikzpicture}
\end{center}
破線は状態方程式の全枝、実線は平衡等温線である。共存区間は一定圧力の直線で結び、区間外では実線と状態方程式が一致する。

\end{document}
